\chapter{Espacios Vectoriales Booleanos y Códigos de Hamming}

\section{El Espacio Vectorial Binario $\mathbb{B}^n$}

En el Capítulo 6b demostramos que, si bien cualquier álgebra de Boole puede ser interpretada como un anillo, el \textbf{único} álgebra de Boole que tiene la estructura rigurosa de un \textbf{Cuerpo Matemático} (es decir, carente de divisores de cero y donde todo elemento no nulo tiene inverso multiplicativo) es el álgebra bivaluada $\mathbb{B}_2 = \{0, 1\}.$ Algebraicamente, este cuerpo es exactamente el cuerpo de Galois $\mathbb{F}_2.$

Dado que la condición insoslayable para construir un espacio vectorial lineal es operar sobre un cuerpo de escalares, deducimos que los únicos espacios vectoriales puramente booleanos que pueden existir deben tener como conjunto base a $\mathbb{F}_2.$

\begin{quote}\textbf{Definicion (El Espacio Vectorial $\mathbb{B}^n$):}\label{espacio_bn}
Se define el espacio vectorial booleano $\mathbb{B}^n$ como el conjunto de todas las $n$-tuplas (vectores de $n$ bits) cuyos elementos pertenecen a $\mathbb{F}_2.$ 
Las dos operaciones que dotan al conjunto de estructura de espacio vectorial son:
\begin{enumerate}
    \item \textbf{Suma vectorial:} Se define como la operación XOR ($\oplus$) aplicada bit a bit entre dos vectores.
    \[ \vec{u} \oplus \vec{v} = (u_1 \oplus v_1, u_2 \oplus v_2, \ldots, u_n \oplus v_n) \]
    \item \textbf{Producto por escalar:} Se define como la operación AND ($\wedge$) entre un escalar booleano $k \in \mathbb{F}_2$ y cada elemento del vector.
    \[ k \wedge \vec{v} = (k \wedge v_1, k \wedge v_2, \ldots, k \wedge v_n) \]
\end{enumerate}
\end{quote}

Al operar sobre $\mathbb{F}_2,$ este espacio vectorial hereda directamente la característica 2 de su cuerpo base. Esto significa que \textbf{todo vector es su propio inverso aditivo}:
\[ \vec{v} \oplus \vec{v} = \vec{0} \]
lo cual simplifica extraordinariamente el cálculo matricial (sumar y restar vectores es exactamente la misma operación).

\section{Métrica: Peso y Distancia de Hamming}

Para que este espacio abstracto tenga utilidad práctica en ingeniería (en concreto, para analizar cómo de parecidos son dos mensajes digitales), necesitamos dotarlo de una métrica que nos permita medir "distancias" entre vectores.

\begin{quote}\textbf{Definicion (Función Distancia (Métrica)):}\label{metrica}
Dado un conjunto $X,$ una \textbf{métrica} o \textbf{función distancia} sobre $X$ es una función $d: X \times X \to \mathbb{R}$ que, para cualesquiera $x, y, z \in X,$ satisface las siguientes cuatro propiedades:
\begin{enumerate}
    \item \textbf{No negatividad:} $d(x, y) \ge 0.$
    \item \textbf{Identidad de los indiscernibles:} $d(x, y) = 0 \iff x = y.$
    \item \textbf{Simetría:} $d(x, y) = d(y, x).$
    \item \textbf{Desigualdad triangular:} $d(x, z) \le d(x, y) + d(y, z).$
\end{enumerate}
\end{quote}

\begin{quote}\textbf{Definicion (Peso de Hamming):}\label{peso_hamming}
El \textbf{peso de Hamming} de un vector $\vec{v} \in \mathbb{B}^n,$ denotado como $w(\vec{v}),$ es el número de componentes no nulas (número de unos) que contiene. En términos formales, si $\vec{v} = (v_1, v_2, \dots, v_n),$ entonces:
\[ w(\vec{v}) = \sum_{i=1}^n v_i \]
donde la suma se entiende en la aritmética ordinaria de $\mathbb{R}.$
\end{quote}

\begin{quote}\textbf{Definicion (Distancia de Hamming):}\label{distancia_hamming}
La \textbf{distancia de Hamming} entre dos vectores $\vec{u}, \vec{v} \in \mathbb{B}^n,$ denotada como $d(\vec{u}, \vec{v}),$ es el número de posiciones en las que difieren. Operacionalmente, coincide con el peso de Hamming de su suma vectorial en $\mathbb{B}^n:$
\[ d(\vec{u}, \vec{v}) = w(\vec{u} \oplus \vec{v}) \]
\end{quote}

\begin{quote}\textbf{Teorema (La Distancia de Hamming es una Métrica):}\label{hamming_metrica}
La función $d(\vec{u}, \vec{v}) = w(\vec{u} \oplus \vec{v})$ cumple todas las propiedades matemáticas de una métrica sobre el espacio $\mathbb{B}^n.$
\end{quote}
\begin{proof}
Debemos verificar el cumplimiento de las cuatro propiedades definitorias para cualesquiera vectores $\vec{u}, \vec{v}, \vec{z} \in \mathbb{B}^n:$
\begin{enumerate}
    \item \textbf{No negatividad:} Dado que el peso de Hamming $w(\vec{x})$ cuenta el número de bits a $1$ en un vector, es una suma de números enteros no negativos, por lo que $w(\vec{u} \oplus \vec{v}) \ge 0.$ Por definición, $d(\vec{u}, \vec{v}) \ge 0.$
    \item \textbf{Identidad de los indiscernibles:}
    \[ d(\vec{u}, \vec{v}) = 0 \iff w(\vec{u} \oplus \vec{v}) = 0 \]
    El peso de un vector es cero si y solo si todos sus componentes son cero, es decir, $\vec{u} \oplus \vec{v} = \vec{0}.$ En el espacio booleano, dos elementos suman cero (vía XOR) si y solo si son idénticos, por tanto $\vec{u} = \vec{v}.$
    \item \textbf{Simetría:}
    Por la conmutatividad intrínseca de la operación XOR en el cuerpo $\mathbb{F}_2:$
    \[ d(\vec{u}, \vec{v}) = w(\vec{u} \oplus \vec{v}) = w(\vec{v} \oplus \vec{u}) = d(\vec{v}, \vec{u}) \]
    \item \textbf{Desigualdad triangular:} Demostraremos que $d(\vec{u}, \vec{z}) \le d(\vec{u}, \vec{v}) + d(\vec{v}, \vec{z})$ mediante dos enfoques distintos.
    
    \textit{Enfoque A (Algebraico a través del peso):}
    Primero, observemos una propiedad fundamental del peso respecto a la suma vectorial: un bit en $\vec{a} \oplus \vec{b}$ solo es $1$ si los bits correspondientes en $\vec{a}$ y $\vec{b}$ son distintos. Por tanto, los $1$s en $\vec{a} \oplus \vec{b}$ provienen como máximo de la unión de los $1$s presentes en $\vec{a}$ y en $\vec{b}.$ Esto implica que $w(\vec{a} \oplus \vec{b}) \le w(\vec{a}) + w(\vec{b}).$
    Si tomamos $\vec{a} = \vec{u} \oplus \vec{v}$ y $\vec{b} = \vec{v} \oplus \vec{z},$ obtenemos:
    \[ w((\vec{u} \oplus \vec{v}) \oplus (\vec{v} \oplus \vec{z})) \le w(\vec{u} \oplus \vec{v}) + w(\vec{v} \oplus \vec{z}) \]
    Aplicando asociatividad y el hecho de que $\vec{v} \oplus \vec{v} = \vec{0},$ el lado izquierdo se simplifica a $w(\vec{u} \oplus \vec{z}).$ Sustituyendo la definición de distancia:
    \[ d(\vec{u}, \vec{z}) \le d(\vec{u}, \vec{v}) + d(\vec{v}, \vec{z}) \]
    
    \textit{Enfoque B (Componente a componente):}
    Analicemos una única coordenada $i$ cualquiera. Las posibles combinaciones de valores $(u_i, v_i, z_i) \in \{0, 1\}^3$ y su contribución a la distancia (donde la contribución a $d(\vec{x}, \vec{y})_i$ es $x_i \oplus y_i$) son:
    \begin{itemize}
        \item Si $u_i = z_i,$ entonces $d(\vec{u}, \vec{z})_i = 0.$ Como la suma $d(\vec{u}, \vec{v})_i + d(\vec{v}, \vec{z})_i$ solo contiene valores no negativos (0 o 1), la desigualdad $0 \le d(\vec{u}, \vec{v})_i + d(\vec{v}, \vec{z})_i$ se cumple trivialmente independientemente del valor de $v_i.$
        \item Si $u_i \ne z_i,$ entonces $d(\vec{u}, \vec{z})_i = 1.$ En este caso, el valor intermedio $v_i$ debe ser igual a $u_i$ o igual a $z_i$ (ya que solo hay dos estados posibles en $\mathbb{F}_2$). 
        \begin{itemize}
            \item Si $v_i = u_i \ne z_i,$ entonces $d(\vec{u}, \vec{v})_i = 0$ y $d(\vec{v}, \vec{z})_i = 1,$ cuya suma es $1.$
            \item Si $v_i = z_i \ne u_i,$ entonces $d(\vec{u}, \vec{v})_i = 1$ y $d(\vec{v}, \vec{z})_i = 0,$ cuya suma es $1.$
        \end{itemize}
        En ambos sub-casos obtenemos $1 \le 1,$ verificándose la desigualdad para la coordenada $i.$
    \end{itemize}
    Dado que la distancia total es la suma de las contribuciones de las $n$ componentes, y la desigualdad se sostiene componente a componente, la desigualdad triangular global queda demostrada.
\end{enumerate}
\end{proof}

\begin{quote}\textbf{Teorema (Invarianza Traslacional de la Distancia):}\label{inv_traslacional}
La distancia de Hamming es invariante bajo traslaciones. Es decir, sumar un mismo vector $\vec{w}$ a dos vectores dados no altera la distancia entre ellos:
\[ d(\vec{u} \oplus \vec{w}, \vec{v} \oplus \vec{w}) = d(\vec{u}, \vec{v}) \]
\end{quote}
\begin{proof}
Partimos de la definición de distancia aplicada a los vectores trasladados:
\[ d(\vec{u} \oplus \vec{w}, \vec{v} \oplus \vec{w}) = w((\vec{u} \oplus \vec{w}) \oplus (\vec{v} \oplus \vec{w})) \]
Por las propiedades de asociatividad y conmutatividad de la operación suma en nuestro espacio vectorial, podemos reordenar los términos:
\[ = w(\vec{u} \oplus \vec{v} \oplus (\vec{w} \oplus \vec{w})) \]
Sabemos que cualquier vector sumado consigo mismo resulta en el vector nulo ($\vec{w} \oplus \vec{w} = \vec{0}$) en un espacio de característica $2.$ Así:
\[ = w(\vec{u} \oplus \vec{v} \oplus \vec{0}) = w(\vec{u} \oplus \vec{v}) = d(\vec{u}, \vec{v}) \]
\end{proof}

\begin{quote}\textbf{Teorema (Acotación de la Distancia en $\mathbb{B}^n$):}\label{cota_hamming}
En el espacio vectorial bivaluado de dimensión finita $\mathbb{B}^n,$ la distancia de Hamming entre cualquier par de vectores está acotada estrictamente superior e inferiormente:
\[ 0 \le d(\vec{u}, \vec{v}) \le n \]
\end{quote}
\begin{proof}
Por la propia definición axiomática de la métrica que hemos comprobado previamente, el límite inferior $d(\vec{u}, \vec{v}) \ge 0$ es trivial (no negatividad).

Para el límite superior, consideremos la definición $d(\vec{u}, \vec{v}) = w(\vec{u} \oplus \vec{v}).$ El peso de Hamming de un vector cuenta el número de posiciones no nulas. Puesto que los vectores en $\mathbb{B}^n$ están formados por exactamente $n$ componentes (es decir, tienen dimensión $n$), el número máximo de componentes que pueden ser distintas de cero (y por tanto, el número máximo de posiciones en las que $\vec{u}$ y $\vec{v}$ difieren) es el número total de componentes del vector, que es $n.$ 

En el caso extremo en el que los vectores difieren en todos los bits (siendo uno el complemento bit a bit del otro, $\vec{v} = \neg \vec{u}$), obtenemos:
\[ d(\vec{u}, \neg \vec{u}) = w(\vec{u} \oplus \neg \vec{u}) = w(\vec{1}) = n \]
Por tanto, la distancia jamás puede exceder la dimensión $n$ del espacio.
\end{proof}

\section{Códigos de Corrección de Errores (Hamming)}

La aplicación más brillante de dotar a las cadenas de bits de una estructura de espacio vectorial sobre $\mathbb{F}_2$ es la invención de los \textbf{códigos correctores de errores}. 

En un canal de comunicación ruidoso, un vector $\vec{v}$ enviado puede sufrir corrupciones (cambios de 0 a 1 o viceversa), recibiéndose un vector diferente $\vec{r}.$ Richard Hamming propuso solucionar esto no usando todo el espacio vectorial $\mathbb{B}^n,$ sino limitando los mensajes válidos a un subespacio vectorial más pequeño y controlado.

\subsection{El Subespacio Código y la Matriz de Paridad}
Un código lineal por bloques de longitud $n$ y dimensión $k$ se define matemáticamente como un \textbf{subespacio vectorial} $C \subset \mathbb{B}^n$ de dimensión $k.$ Todo subespacio vectorial puede definirse como el núcleo (kernel) de una transformación lineal, representada por una matriz llamada \textbf{Matriz de Paridad ($H$)} de dimensiones $(n-k) \times n.$

\begin{quote}\textbf{Teorema (Validación del Síndrome):}\label{sindrome_hamming}
Un vector recibido $\vec{r}$ es una palabra código válida (pertenece al subespacio código $C$) si y solo si su producto por la matriz de paridad $H$ (utilizando aritmética en $\mathbb{F}_2$) da como resultado el vector nulo. A este resultado se le denomina \textbf{síndrome} ($\vec{s}$).
\[ \vec{s} = H \cdot \vec{r}^T \]
Si $\vec{s} = \vec{0},$ el vector pertenece al subespacio (no hay errores detectados). Si $\vec{s} \ne \vec{0},$ el vector ha salido del subespacio, lo que indica que se ha corrompido durante la transmisión.
\end{quote}

Dado que operar matrices sobre $\mathbb{F}_2$ requiere únicamente puertas lógicas XOR y AND, el cálculo del síndrome $\vec{s} = H \cdot \vec{r}^T$ se puede implementar directamente en hardware digital con una eficiencia extrema, constituyendo la base de los modernos sistemas de memoria ECC (Error-Correcting Code).
